\subsection{测度与紧集的可加集合函数}

定理 1. --- 设 T 是拓扑空间，I 是从 \(\mathfrak{K}(\mathrm{T})\) 到 \(\mathbf{R}_{+}\) 的映射。要使存在 T 上的测度 \(\mu\)，使得 \(\mathrm{I}(\mathrm{K}) = \mu^{\bullet}(\mathrm{K})\) 对所有 \(\mathrm{K} \in \mathfrak{K}(\mathrm{T})\) 成立，必要且充分的条件是 I 满足以下条件：

\begin{enumerate}
\def\labelenumi{\arabic{enumi})}
\item
  如果 K 和 L 是 T 的紧子集且 \(K \subset L\)，则 \(I(K) \leqslant I(L)\)（“I 递增”）。
\item
  如果 K 和 L 是 T 的紧子集，则 I(K ∪ L) ≤ I(K) + I(L)。
\item
  如果 K 和 L 是 T 的不交紧子集，则 \(\mathrm{I}(\mathrm{K}\cup\mathrm{L})=\mathrm{I}(\mathrm{K})+\mathrm{I}(\mathrm{L})\)（“I 可加”）。
\item
  对于 T 的每个紧子集的递减有向族 \((\mathrm{K}_{\alpha})_{\alpha \in \mathrm{A}}\)，都有 \(\mathrm{I}\big(\bigcap_{\alpha \in \mathrm{A}}\mathrm{K}_{\alpha}\big) = \inf_{\alpha \in \mathrm{A}}\mathrm{I}(\mathrm{K}_{\alpha})\)。
\item
  对于每个 \(x \in \mathbf{T}\)，存在一个邻域 \(\mathbf{V}\) 的 \(x\)，使得
\end{enumerate}

\[
\sup_{\substack{\mathbf{K}\in \mathfrak{K}(\mathbf{T})\\ \mathbf{K}\subset \mathbf{V}}}\operatorname{I}(\mathbf{K}) <   + \infty
\]

（“I 局部有界”）。

测度 \(\mu\) 于是唯一。

\(\mu\) 的唯一性由 §1, No.~2 的 Cor. of Prop. 2 得出。上述条件是必要的：前三个条件显然，最后一个条件来自 \(\mu^{\bullet}\) 是局部有界负荷这一事实，而条件 4) 由 §1, No.~6 的 Cor. of Prop. 5 得出。

为证明这些条件是充分的，我们先处理 \(\mathbf{T}\) 为紧空间的情形。

引理 1. --- 假设 T 是紧的，并置 \(l = \mathrm{I}(\mathrm{T})\)。对于每个 \(\mathrm{A} \subset \mathrm{T}\)，置

\[
\mathrm{J}(\mathrm{A}) = \sup_{\substack{\mathrm{K}\in \mathfrak{K}(\mathrm{T})\\ \mathrm{K}\subset \mathrm{A}}}\mathrm{I}(\mathrm{K})\tag{1}
\]

并令 \(\Phi\) 为满足 \(A \subset T\) 且 \(J(A) + J(\mathbb{C}A) = l\) 的集合。于是 \(\Phi\) 是包含 \(\mathfrak{K}(T)\) 的集环，且函数 \(J\) 在 \(\Phi\) 上递增并可加。

显然，\(J\) 是一个递增的集合函数，延拓 \(I\)，并且 \(J(A) + J(\mathbb{C}A) \leqslant l\) 对每个 \(A \subset T\) 都成立。

令 K 和 S 是 T 中的两个紧集；我们首先要证明

\[
\mathrm{J} (\mathrm{K} \cap \mathrm{S}) + \mathrm{J} (\mathbb {C} \mathrm{K} \cap \mathrm{S}) = \mathrm{J} (\mathrm{S}).\tag{2}
\]

考察 I 在 \(\mathfrak{K}(\mathrm{S})\) 上的限制和 J 在 \(\mathfrak{P}(\mathrm{S})\) 上的限制，立即可归结为 \(\mathrm{S} = \mathrm{T}\) 的情形。由于 T 是正规空间，K 是其紧邻域的递减有向族的交，而条件 4) 蕴含：对于每个 \(\varepsilon > 0\)，存在 K 的紧邻域 H，使得 \(\mathrm{I}(\mathrm{H}) \leqslant \mathrm{I}(\mathrm{K}) + \varepsilon\)。令 L 为 T - H 的闭包；L 是紧的，\(\mathrm{L} \cap \mathrm{K} = \varnothing\) 且 \(\mathrm{H} \cup \mathrm{L} = \mathrm{T}\)，因此 \(l = \mathrm{I}(\mathrm{H} \cup \mathrm{L}) \leqslant \mathrm{I}(\mathrm{H}) + \mathrm{I}(\mathrm{L}) \leqslant \mathrm{I}(\mathrm{K}) + \mathrm{I}(\mathrm{L}) + \varepsilon\)（条件 2)），从而关系 \(\mathrm{J}(\mathrm{K}) + \mathrm{J}(\mathbb{C}\mathrm{K}) \geqslant \mathrm{I}(\mathrm{K}) + \mathrm{I}(\mathrm{L}) \geqslant l - \varepsilon\)。由于 \(\varepsilon\) 任意，\(\mathrm{J}(\mathrm{K}) + \mathrm{J}(\mathbb{C}\mathrm{K}) = l\)。这证明了公式 (2)，也证明了包含关系 \(\mathfrak{K}(\mathrm{T}) \subset \Phi\)。

现在证明 \(\Phi\) 是一个集环。由于 \(\Phi\) 显然对取补稳定，只需证明：如果 \(A_{1}\) 和 \(A_{2}\) 是 \(\Phi\) 的元素，则 \(A_{1} \cup A_{2} \in \Phi\)，或者说，

\[
\mathrm{J} \left(\mathrm{A} _ {1} \cup \mathrm{A} _ {2}\right) + \mathrm{J} \left(\mathbf {C} \left(\mathrm{A} _ {1} \cup \mathrm{A} _ {2}\right)\right) \geqslant l.\tag{3}
\]

令 \(\varepsilon\) 是一个 \(>0\) 的数，并且对于 \(i = 1,2\)，令 \(\mathbf{K}_i\) 是包含于 \(\mathbf{A}_i\) 的紧集，令 \(\mathbf{L}_i\) 是包含于 \(\mathbf{C}\mathbf{A}_i\) 的紧集，使得

\[
\mathrm{I} \left(\mathrm{K} _ {i}\right) \geqslant \mathrm{J} \left(\mathrm{A} _ {i}\right) - \varepsilon , \quad \mathrm{I} \left(\mathrm{L} _ {i}\right) \geqslant \mathrm{J} \left(\mathbb {C} \mathrm{A} _ {i}\right) - \varepsilon .
\]

置 \(\mathbf{M}_1 = \mathbf{K}_1\cup \mathbf{L}_1\)；关系 \(l = J(\mathbf{M}_1) + \mathbf{J}(\mathbf{C}\mathbf{M}_1)\) 以及

\[
\mathrm{J} \left(\mathrm{M} _ {1}\right) = \mathrm{I} \left(\mathrm{K} _ {1}\right) + \mathrm{I} \left(\mathrm{L} _ {1}\right) \geqslant \mathrm{J} \left(\mathrm{A} _ {1}\right) + \mathrm{J} \left(\mathbf {C A} _ {1}\right) - 2 \varepsilon = l - 2 \varepsilon
\]

蕴含 \(\mathrm{J}(\mathbb{C}\mathrm{M}_1) \leqslant 2\varepsilon\)。于是若 \(\mathrm{S}\) 是 \(\mathrm{T}\) 的紧子集，关系 (2)（应用于 \(\mathrm{K} = \mathrm{M}_1\)）蕴含 \(\mathrm{J}(\mathrm{S}) \leqslant \mathrm{J}(\mathrm{M}_1 \cap \mathrm{S}) + 2\varepsilon\)，从而

\[
\mathrm{J} (\mathrm{S}) \leqslant \mathrm{J} (\mathrm{K} _ {1} \cap \mathrm{S}) + \mathrm{J} (\mathrm{L} _ {1} \cap \mathrm{S}) + 2 \varepsilon .
\]

将取 \(S = K_{2}\) 和 \(S = L_{2}\) 所得的不等式相加，并考虑不等式 \(\mathrm{J}(K_{2}) + \mathrm{J}(L_{2}) \geqslant l - 2\varepsilon\)，以及 \(K_{1} \cap K_{2}\)、\(L_{1} \cap K_{2}\) 和 \(K_{1} \cap L_{2}\) 是包含于 \(A_{1} \cup A_{2}\) 的三个不交紧集这一事实。以 C 表示这三个紧集的并，则

\[
\begin{array}{r l} & {l - 2 \varepsilon \leqslant \mathrm{J} (\mathrm{K} _ {2}) + \mathrm{J} (\mathrm{L} _ {2}) \leqslant \mathrm{J} (\mathrm{C}) + \mathrm{J} (\mathrm{L} _ {1} \cap \mathrm{L} _ {2}) + 4 \varepsilon} \\ & {\qquad \leqslant \mathrm{J} (\mathrm{A} _ {1} \cup \mathrm{A} _ {2}) + \mathrm{J} \big (\mathbb {C} (\mathrm{A} _ {1} \cup \mathrm{A} _ {2}) \big) + 4 \varepsilon ,} \end{array}
\]

由此立即得到所需公式 (3)，因为 \(\varepsilon\) 是任意的。既然这已经建立，前面的不等式蕴含 \(\mathrm{J}(\mathbf{C}) \geqslant \mathrm{J}(\mathbf{A}_1 \cup \mathbf{A}_2) - 6\varepsilon\)；如果 \(\mathbf{A}_1\) 和 \(\mathbf{A}_2\) 不交，则 \(\mathbf{C}\) 是 \(\mathbf{K}_1 \cap \mathbf{L}_2 \subset \mathbf{A}_1\) 与 \(\mathbf{K}_2 \cap \mathbf{L}_1 \subset \mathbf{A}_2\) 的并，由此可推出 \(\mathrm{J}(\mathbf{A}_1 \cup \mathbf{A}_2) \leqslant \mathrm{J}(\mathbf{A}_1) + \mathrm{J}(\mathbf{A}_2)\)。反向不等式显然成立，故 \(\mathbf{J}\) 确实在 \(\Phi\) 上可加，引理得证。

我们完成 T 为紧空间时的定理证明。令 \(\mathcal{E}(\Phi)\) 为 T 上的 \(\Phi\)-阶梯函数所成的向量空间，赋予一致收敛拓扑（Ch. IV, §4, No.~9, Def. 4）；仍以 J 表示与可加函数 J 关联的 \(\mathcal{E}(\Phi)\) 上的正线性形式（loc. cit., Prop. 18）。由于 \(J(T) = l\)，J 连续且范数为 \(l\)。令 \(\mathcal{H}\) 为 \(\mathcal{E}(\Phi)\) 在一致收敛拓扑下的闭包；立即可验证，J 可由连续性延拓为 \(\mathcal{H}\) 上的正线性形式，仍记为 J。由于 \(\mathcal{H}\) 包含 \(\mathcal{C}(T)\)（loc. cit., No.~10, Prop. 19），J 在 \(\mathcal{C}(T)\) 上的限制是正测度 \(\mu\)。还需证明 \(\mu^{\bullet}(K) = I(K)\)，对于 T 的每个紧子集 K。现在有 \(\mu^{\bullet}(K) = \inf_{f \in S_K} \mu^{\bullet}(f)\)，其中 \(S_K\) 表示 \(\mathcal{C}(T)\) 中满足 \(\geqslant \varphi_K\) 的元素所成的集合（§1, No.~6, Prop. 5）。由于 \(J(f) = \mu^{\bullet}(f)\) 对于 \(f \in \mathcal{C}(T)\) 成立，显然只需证明 \(J(K) \geqslant \inf_{f \in S_K} J(f)\)。如引理 1 的证明中那样，令 H 是 K 的紧邻域，使得 \(J(H) \leqslant J(K) + \varepsilon\)，并令 \(f\) 是 T 上取值介于 0 和 1 之间、在 K 上等于 1 且在 H 外为 0 的连续函数（GT, IX, §4, No.~1, Prop. 1）。于是

\[
\mathrm{J} (f) \leqslant \mathrm{J} (\mathrm{H}) \leqslant \mathrm{J} (\mathrm{K}) + \varepsilon ;
\]

\(\varepsilon\) 任意，所需不等式得证，定理在 \(\mathbf{T}\) 为紧空间时也就得证。

现在转到一般情形。对于 T 中的每个紧集 L，令 \(I_{L}\) 为 I 在 \(\mathfrak{K}(L)\) 上的限制。根据刚处理的特殊情形，存在 L 上唯一的测度 \(\mu_{L}\)，使得 \(\mu_{\mathrm{L}}(\mathrm{K}) = \mathrm{I}_{\mathrm{L}}(\mathrm{K})\) 对每个 \(K \subset L\) 的紧集成立。再令 \(L'\) 是包含于 L 的紧集；有 \((\mu_{\mathrm{L}})_{\mathrm{L}'}^{\bullet}(\mathrm{K}) = \mu_{\mathrm{L}}^{\bullet}(\mathrm{K}) = \mu_{\mathrm{L}'}^{\bullet}(\mathrm{K})\) 对每个 \(K \subset L'\) 的紧集成立，因此 \(\mu_{\mathrm{L}'} = (\mu_{\mathrm{L}})_{\mathrm{L}'}\)；映射 \(\mu : L \mapsto \mu_{L}\) 是一个预测度。条件 5) 表明 \(\mu\) 是测度，而关系 \(\mathrm{I}(\mathrm{K}) = \mu^{\bullet}(\mathrm{K})\) 对所有紧集 \(\mathbf{K} \subset \mathbf{T}\) 显然成立。

注记。--- 1) 在定理 1 的叙述中，条件 4) 可以由下述条件（“右连续性”）替代：

4') 对于每个 \(\mathbf{K} \in \mathfrak{K}(\mathbf{T})\) 和每个 \(\varepsilon > 0\)，存在包含 \(\mathbf{U}\) 的开集 \(\mathbf{K}\)，使得对于每个紧集 \(\mathrm{I}(\mathrm{H}) \leqslant \mathrm{I}(\mathrm{K}) + \varepsilon\)，都有 \(\mathrm{H} \subset \mathbf{U}\)。

事实上，若 \(\mu\) 是测度，则函数 \(I: K \mapsto \mu^{\bullet}(K)\) 满足 \(4'\)（§1, No.~9, Prop. 13）。反之，假设 I 满足 1) 和 \(4'\)；我们来证明 I 于是满足 4)。沿用定理 1 叙述中的记号，取一个 \(\varepsilon > 0\) 以及包含紧集 \(K = \bigcap_{\alpha \in A} K_{\alpha}\) 的开集 U，并使得 \(4'\) 成立。

于是存在指标 \(\beta \in A\)，使得 \(K_{\beta} \subset U\)，这蕴含

\[
\inf _ {\alpha \in \mathbf {A}} \mathrm{I} (\mathrm{K} _ {\alpha}) \leqslant \mathrm{I} (\mathrm{K} _ {\beta}) \leqslant \mathrm{I} (\mathrm{K}) + \varepsilon
\]

且 4) 确实得到验证。

\begin{enumerate}
\def\labelenumi{\arabic{enumi})}
\setcounter{enumi}{1}
\tightlist
\item
  条件 2) 和 3) 的集合可以在定理 1 的叙述中由下述条件替代：
\end{enumerate}

如果 K 和 L 是 T 的紧子集，则

\[
\mathrm{I} (\mathrm{K} \cup \mathrm{L}) + \mathrm{I} (\mathrm{K} \cap \mathrm{L}) = \mathrm{I} (\mathrm{K}) + \mathrm{I} (\mathrm{L}).
\]

事实上，这个条件蕴含 2) 和 3)，另一方面

\[
\mu^ {\bullet} (\mathrm{K} \cup \mathrm{L}) = \mu^ {\bullet} (\mathrm{K} \cap \mathrm{L}) = \mu^ {\bullet} (\mathrm{K}) + \mu^ {\bullet} (\mathrm{L})
\]

对于每个测度 \(\mu\)，这是由特征函数之间的关系 \(\varphi_{\mathrm{K}\cup\mathrm{L}} + \varphi_{\mathrm{K}\cap\mathrm{L}} = \varphi_{\mathrm{K}} + \varphi_{\mathrm{L}}\) 得到的。

